\section{失败模式：优化 reward 不等于优化真实价值}

回到整个 post-training pipeline 的边界，无论 PPO 还是 DPO，都在优化一个代理信号。随着优化增强，模型可能越来越善于利用 reward model，而不是越来越符合真实目标；因此验收必须使用训练信号之外的独立测量。

\singleslide{slide-62.jpg}{RLHF 的两类总风险：reward overoptimization，以及输出分布的 mode collapse / entropy loss。}{62}

Slide 62 把后果分成“目标错位”和“分布退化”。前者关注模型是否利用 reward 的漏洞，后者关注即使平均 reward 上升，policy 是否失去多样性、校准与少数合理答案。两类问题可以同时发生，因此只看 reward curve 或单次 greedy 输出都不足以验收 post-training。

\subsection{Reward overoptimization}

本节先看最直接的 Goodhart 曲线：训练 reward 持续上升，并不保证 held-out human preference 同步上升。早期 update 修复明显缺陷，后期 policy 则可能进入 reward model 未覆盖的分布，利用长度、格式、特定短语或评测器盲点获得高分。

\begin{figure}[H]
\centering
\includegraphics[width=0.84\textwidth]{slide-63.jpg}
\caption{继续优化 proxy reward 后，真实偏好性能可能先升后降。}
\figsource{CS336 Spring 2026 Lecture 15，Slide 63。}
\end{figure}

\begin{knowledgebox}{读图：最优点为什么出现在中间}
早期优化提高有用行为；继续训练后，policy 进入 reward model 的 distribution shift 区域，学会 verbosity、格式或其他 shortcut。曲线转折意味着 validation reward、held-out human eval 和 policy KL 必须同时监控。
\end{knowledgebox}

\subsection{Mode collapse 与 calibration loss}

进一步看分布形状，即使平均 preference score 提高，policy 也可能把多个合理答案压成单一模板。概率质量集中后，模型的 entropy、minority viewpoint 和不确定性表达都会下降；这类退化无法从 top-1 win rate 单独发现。

\begin{figure}[H]
\centering
\includegraphics[width=0.84\textwidth]{slide-64.jpg}
\caption{RLHF 可能压缩输出分布，降低多样性与概率校准。}
\figsource{CS336 Spring 2026 Lecture 15，Slide 64。}
\end{figure}

Base LM 是 probabilistic model，能够表达多个合理续写；强偏好优化可能把分布压向单一高 reward style，使 entropy 降低、minority answers 消失、uncertainty 不再反映真实正确率。对 brainstorming、creative writing 和 ambiguous questions，这种 collapse 尤其明显。

老师在课程结束前特别补充，entropy loss 与 mode collapse 虽然只剩很短时间讨论，却不应被当成次要问题。它们改变的不是某个 benchmark 分数，而是模型作为概率分布的性质：采样多次仍只得到同一种话术，置信表达也不再与事实正确率对应。

\begin{importantbox}{Post-training 的正确验收方式}
不要只看一个 aggregate preference score。至少同时检查 task correctness、safety slices、length、calibration、diversity、KL drift、rare domains、tool reliability 和 adversarial prompts。
\end{importantbox}

\begin{longtable}{L{0.2\textwidth}L{0.32\textwidth}L{0.38\textwidth}}
\toprule
监控信号 & 发现什么 & 需要的独立对照 \\
\midrule
Train/reward-model score & 优化器是否在提高代理目标 & 新版本 judge 与 human blind eval \\
Policy KL & 离 reference 多远 & 分 domain、长度和语言切片 \\
Response length/style & 是否靠 verbosity 获胜 & length-matched pairwise eval \\
Entropy/diversity & 是否出现 mode collapse & 多次采样、distinct strategy/answer count \\
Calibration & 置信度是否仍对应正确率 & 可验证问题上的 reliability diagram \\
Tool success & 格式正确是否等于任务完成 & 执行结果、side effect 与恢复能力 \\
Safety slices & 平均安全是否掩盖长尾 & 多语言、多轮、工具权限和 adversarial set \\
\bottomrule
\caption{Post-training 失败模式监控表。}
\end{longtable}

\begin{knowledgebox}{Nota de clase}
讲者最后强调，RLHF 文献中的胜负常高度依赖实验配置。算法名称不是结果的充分解释；数据、base model、reward quality、sampling 和 evaluation 同样重要。
\end{knowledgebox}

\subsection{本章小结}

Reward optimization 的根本风险是 Goodhart's law：代理指标一旦成为目标，就更容易被利用。稳定 post-training 需要限制分布漂移、持续 fresh evaluation，并保留对模型多样性和校准的观察。

\singleslide{slide-65.jpg}{全讲回顾：RLHF 的瓶颈同时存在于数据收集、优化系统和代理目标过度优化。}{65}

最终 recap 把课程压缩成三条，但应按因果顺序理解：先有困难且带偏差的 feedback collection，随后 PPO/DPO 把这些信号写回 policy，最后更强的优化又可能放大 reward misspecification。没有独立 human/verifier evaluation、分布切片和停止规则，任何单一训练分数都不能证明模型真正变好。

